% =====================================================================
%  第 6 章 Gromov–Witten 不变量与镜像对称 · §6.3 Frobenius 流形
%  ch06b —— 页段：PDF 286–303（原书 p.276–293）
% =====================================================================

% [原文 p.276]
% 注：本节标题前的两句是 §6.2 末尾（原书 p.276 上半）。上一句“……[3] 是优秀的教科书。”已由 ch06a2 译出，此处从下一句接续。
把 operad 用于弦理论是比较新的事。参见 [198], [238] 等。共形场论中使用的顶点算子代数（vertex operator algebra\footnote{参见 [171] 等。}）等，也与 operad 关系很深。它与 §6.4、§6.5、§6.6 中用到的 $L_\infty$ 代数、$A_\infty$ 代数也密切相关。

\section{Frobenius 流形}

\takeaway{本节补上 §6.1 讲 $\mathcal{SJ}(M)$ 时略去的一点：$\mathcal{SJ}(M)$ 应是“扩展模空间”在 $M$ 处的切空间。技术上的落实是把 Gromov–Witten 势写成模空间上的形式幂级数，再用它族化出以 $H^*(M;\C)$ 为参数的 Frobenius 代数结构。}

§6.1 中关于 $\mathcal{SJ}(M)$ 的说明漏掉了一个要点，此处补述。

\textbf{观察 6.73}\quad $X$ 型理论情形的 $\mathcal{SJ}(M)$，是构造 $X$ 型理论时赋予 $M$ 的几何结构的扩展模空间在 $M$ 处的切空间。\hfill$\square$

也就是说，$A$ 模型时辛结构的扩展模空间（extended moduli space）的切空间是 $\mathcal{SJ}^A(M)$，$B$ 模型时复结构的扩展模空间的切空间是 $\mathcal{SJ}^B(M)$。这里用了“扩展模空间”一词，其含义在本章中会逐渐清楚。（复结构的扩展形变空间这一想法，我想始于 Ran [312]。）

观察 6.73 在物理上的说明大致如下。

% [Tutor Note: 观察 6.73 是全节的出发点：$\mathcal{SJ}(M)$ 不是“模空间”，而是模空间的切空间。$A$ 模型的“扩展”模空间把 $H^2(M;\R)$ 方向之外的形变（即 $\bigoplus_{n\ne1}H^{2n}$）也算进来，这一点在数学上就体现为下文以 $H^*(M;\C)$ 为参数的 Frobenius 代数族。]

要构造场论，需考虑 Lagrange 泛函 $\mathcal{L}\colon X\to\R$，并考察其极小值附近的“量子涨落”。$\mathcal{L}$ 在临界点 $p$ 附近可用二次函数、即 Hessian 矩阵 $\operatorname{Hess}_p\mathcal{L}$ 近似。用这个二次的 Lagrange 泛函来考虑，就得到没有相互作用（＝非线性项）的理论，即自由场。自由场中表示粒子的线性空间（Hilbert 空间）是 $p$ 处的切空间 $T_pX$。各个粒子的质量是 $\operatorname{Hess}_p\mathcal{L}$ 的特征值。

这里写的 $\mathcal{SJ}(M)$ 之类，并不是上述 Hilbert 空间 $T_pX$。$\mathcal{SJ}(M)$ 等所确定的场论，是把 $T_pX$ 所表示的粒子中重粒子的效应“积分”掉、只保留轻粒子及其相互作用的部分

% [原文 p.277]
之后得到的低能有效场论（low energy effective field theory）\footnote{这一点稍作说明。$\mathfrak{m}_k\colon \mathcal{L}(X)^{\otimes k}\to\mathcal{L}(X)$ 决定 $\mathcal{L}(X)$ 所表示的 $k$ 个粒子 $p_1,\dots,p_k$ 相互作用而成为第 $k+1$ 个粒子 $p_{k+1}$ 的概率 $\langle\mathfrak{m}_k(p_1,\dots,p_k),p_{k+1}\rangle$。这一反应有时直接发生，也包括中途发生若干次反应、先产生更“重”的粒子 $q_i$、它再引起别的反应而继续、最终剩下 $p_{k+1}$ 的情形。这些例如（取 $k=3$）可写成 $\sum_i\langle\mathfrak{m}_2'(p_1,p_2),q_i\rangle\langle\mathfrak{m}_2'(q_i,p_3),p_4\rangle$ 之类的项们的（无限）和。这里 $\mathfrak{m}_k'$ 是 $T_pX^{\otimes k}\to T_pX$ 的映射，$q_i$ 是 $\mathcal{SJ}(X)$ 在 $TX$ 中的正交补空间的基。把这个求和的操作表作“积分”。（以上操作是称为重整化（renormalization）的操作之一种。）}。

因此，$\mathcal{SJ}(M)$ 是 $\operatorname{Hess}_p\mathcal{L}$ 的 $0$ 特征空间。（$T_pX$ 是无限维的，$\mathcal{SJ}(M)$ 是有限维的。）另一方面，真空（暂且）由 $\mathcal{L}$ 的临界点 $p$ 决定而确定。这样的 $p$ 并不唯一。特别是当 Hessian 不可逆时，给出 $\mathcal{L}$ 极小值的点集具有正维数。$\mathcal{L}$ 的临界点全体称为\textbf{真空模空间}（moduli space of vacuum）。由以上所述，$\mathcal{SJ}(M)$ 是真空模空间的切空间，这一点是显然的。

\begin{remark}{6.74}
  以上的说明有几处不准确。首先，几乎所有重要的情形里，$\mathcal{L}$ 关于无限维的群（规范变换群）对称，真空模空间是把 $\mathcal{L}$ 的临界点集合用无限维的群（规范变换群）作商所得的商空间。$\mathcal{SJ}(M)$ 也一样，需要沿规范变换群的轨道方向作商。

  这一补充也有两点不准确。第 1，量子化后的系统的对称性与原来的 $\mathcal{L}$ 的对称性一般不同。（注 4.96 中说明过中心扩张是必要的。对称性破缺的机制还有量子异常（anomaly）等。）

  第 2，把 $\mathcal{L}$ 的临界点集合用规范变换群作商所得的商空间，是经典模空间（classical moduli space），并不是量子化后的场论的真空模空间。也就是说，经典模空间会被各种量子效应量子形变。“经典模空间被量子形变”是什么意思，本章中会讲两种情形。数学上的一般理论还没有。“经典模空间的量子形变”出现并被漂亮地算出的是 Seiberg–Witten 理论 [329]，那里 $u$ 平面 $\C/\{\pm1\}$ 被量子形变
\end{remark}

% [原文 p.278]
了。

以上的说明并不充分，把物理背景讲清楚超出笔者的能力。这里只叙述了理解后面构造之背景所需的最小限度。

下面在 $A$ 模型的情形给观察 6.73 一个数学内容。

先作几点说明，使后面的叙述顺畅。第 1，在 $\mathfrak{m}_k^{A,\mathrm{fm}}$ 的定义 (6.15) 中引入了形式变量 $T$，考虑其形式幂级数。即把 $\mathfrak{m}_k^{A,\mathrm{fm}}$ 看作 $T=0$ 附近 $T$ 的“函数”。把 $\omega$ 换成 $\C\omega$、把 $T$ 换成 $T^{-c}$，(6.15) 仍不变。因此，当 $\omega$ 也同时变动时，$T$ 可以不作变量而视为常数\footnote{这一点我想与注 6.13 所述的事实，即弦理论中耦合常数是 dilaton 的期望值，有关。}。由后面的命题 6.84 可知，取 $T=e^{-1}$ 是自然的。此时 $\mathfrak{m}_k^{A,\mathrm{fm}}$ 可看作 $\omega$ 的同调类 $\in H^2(M;\R)$ 的函数。本来理应如此，但用形式幂级数而非收敛幂级数来考虑时，把 $\omega$ 当作变量处理起来（出于完全非本质的理由）很麻烦。因此，以后在形式幂级数的范围内讨论时，把 $T$ 保留下来。

第 2，刚才写了“$[\omega]\in H^2(M;\R)$ 的函数”，这一点要修改订正。即取闭微分 $2$ 形式 $B$，以 $\tilde\omega=\omega+2\pi\sqrt{-1}B$ 的同调类 $\in H^2(M;\C)$ 为变量。当然，此时式 (6.15) 中应把 $T^{\beta\cap\omega}$ 换成 $T^{\beta\cap\tilde\omega}$。由于 $T=e^{-1}$，令
\begin{equation}
  \mathfrak{m}_k^{A,\mathrm{fm}}
  =\sum_\beta\mathfrak{m}_{k,\beta}^A\,
  e^{-2\pi\sqrt{-1}B\cap\beta}T^{\beta\cap\omega_M}\otimes u^{\beta\cap c_1(M)} . \tag{6.19}
\end{equation}
(6.19) 的右边是 $\Lambda_{\mathrm{nov},\C}=\Lambda_{\mathrm{nov}}\otimes\C$ 系数的同调的元素。$\beta$ 是整系数的同调类，故 (6.19) 在把 $[B]$ 平移 $H^2(M;\Z)$ 的元素时不变。也就是说，变量（同调类）$[\omega+2\pi\sqrt{-1}B]$ 在 $H^2(M;\C)/2\pi\sqrt{-1}H^2(M;\Z)$ 中变动。称 $B$ 为\textbf{B 场}（$B$ field）。

\begin{remark}{6.75}
  必须考虑 B 场的理由之一如下。注

  % [原文 p.279]
  6.22 中蒙混过去了，其实应当是 $\mathcal{H}_A^*(M)\otimes\C=\bigoplus_{m+\ell=*}H_\delta^{m,\ell}(M)$。即不取 $\mathcal{H}_A^*(M)$ 的复化，它就不与 $\mathcal{H}_B^*(M)$ 一致。
\end{remark}

那么，回到 $A$ 模型情形的观察 6.73。它的一个数学表述是：存在以同调群 $H^*(M;\C)$ 为切空间、以模空间为参数集的 Frobenius 代数族。定理 2.15 意味着，局部地看，辛流形的模空间与 $H^2(M;\R)$ 相等。$\mathcal{H}_A^*(M)$ 是各种次数（不限于 $2$ 次）的同调群之和。因此，朝 $H^2(M;\R)$ 以外的方向的形变，也必须看作某种意义上的“辛结构的形变”。所谓扩展模空间，首先就是它的一个例子。

那么，要作出以 $\mathcal{H}_A^*(M)$ 为参数的 Frobenius 结构族，就要用 Gromov–Witten 势。由于在形式幂级数的范围内讨论，需要稍作语言准备。为了以后也能用，稍微推广一下、先引入记号。设 $R$ 是有单位元的交换环。

\begin{definition}{6.76}
  对分次 $R$ 模 $C_*$，令 $BC=\sum_{k=0}^\infty C^{\otimes k}$。把 $C_{i_1}\otimes\cdots\otimes C_{i_k}$ 的元素的次数定为 $i_1+\cdots+i_k$，$BC$ 便是分次模。用下式定义积，则 $BC$ 成为 $R$ 代数：
  \[
    (x_1\otimes\cdots\otimes x_k)\cdot(y_1\otimes\cdots\otimes y_\ell)
    =x_1\otimes\cdots\otimes x_k\otimes y_1\otimes\cdots\otimes y_\ell .
  \]
  用 $x\otimes y-(-1)^{\deg x\deg y}y\otimes x$ 生成的理想作 $BC$ 的商，记作 $EC$。$EC$ 是分次交换的 $R$ 代数。
\end{definition}

把 $C$ 的偶数次部分与奇数次部分分别记作 $C_{\mathrm{even}},C_{\mathrm{odd}}$，则 $EC$ 是 $C_{\mathrm{even}}$ 上的对称代数（多项式环）与 $C_{\mathrm{odd}}$ 上的外代数之张量积。

% [Tutor Note: 定义 6.76 的 $BC$ 是张量代数、$EC$ 是其分次交换化，即 $EC\cong\operatorname{Sym}(C_{\mathrm{even}})\otimes\Lambda(C_{\mathrm{odd}})$。用 $E$ 这个字母是因为它与“operad 的自由代数”构造同型：$E(V)$ 正是把 $V$ 的元素当作 $k$ 元运算（$k$ 任意）所生成的自由分次交换代数，故本节与 §6.2 的 operad 语言相通。]

当 $R$ 是 Novikov 环时，把 $BC,EC$ 完备化。考虑 $\Lambda_{\mathrm{nov}}$ 的元素 $x=\sum_i x_iT^{\lambda_i}\otimes u^{k_i}$。这里假定 $i\ne j$ 时 $(\lambda_i,k_i)\ne(\lambda_j,k_j)$，且 $x_i$ 均不为 $0$。令 $|x|=\inf\lambda_i$。

\begin{lemma*}{6.77}
  关于 $\exp(-|\cdot|)$，$\Lambda_{\mathrm{nov}}$ 是完备赋值环。
\end{lemma*}

\begin{proof}
  $|\cdot|$ 是赋值这一条件，即 $|xy|\ge|x|\,|y|$、$|v+w|\ge\max\{|x|,|y|\}$ 由 $|\cdot|$ 的定义可知。距离 $\operatorname{dist}(x,y)=\exp(-|y-x|)$ 的完备

  % [原文 p.280]
  性由 Novikov 环的定义显然。
\end{proof}

\begin{definition}{6.78}
  所谓带滤过的 $\Lambda_{\mathrm{nov}}$ 模（filtered $\Lambda_{\mathrm{nov}}$ module），是指分次 $\Lambda_{\mathrm{nov}}$ 模 $C_*$，并存在 $|\cdot|\colon C_*\to\R_{\ge0}$，使对 $v,w\in C_*$、$x\in\Lambda_{\mathrm{nov}}$ 具有下述性质。
  \begin{enumerate}[label=（\arabic*）,leftmargin=2.4em]
    \item $|xv|\ge|x|\,|v|$。
    \item $|v+w|\ge\max\{|v|,|w|\}$。
    \item $\operatorname{dist}(v,w)=\exp(-|w-v|)$ 给出 $C_*$ 上的完备距离。
  \end{enumerate}
\end{definition}

\begin{remark}{6.79}
  带滤过的 $\Lambda_{\mathrm{nov}}$ 模，除次数之外还确定了滤过结构。即令 $\mathcal{F}_\lambda C_*=\{x\in C_*\mid|x|\ge\lambda\}$，则 $\mathcal{F}_\lambda C_*$ 是 $C_*$ 的部分 $\Lambda_{\mathrm{nov}}$ 模，且当 $\lambda<\lambda'$ 时有 $\mathcal{F}_\lambda C_*\supset\mathcal{F}_{\lambda'}C_*$。与第 5 章一样，次数与 Chern 类相关，滤过与能量相关；在一般辛流形上这两者互不相关。
\end{remark}

\begin{definition}{6.80}
  对带滤过的 $\Lambda_{\mathrm{nov}}$ 模，考虑定义 6.76 的 $BC,EC$。令 $|x_1\otimes\cdots\otimes x_k|=|x_1|+\cdots+|x_k|$，则定义 6.78 的 (1),(2) 得到满足。关于 $\operatorname{dist}(x,y)=\exp(-|y-x|)$ 的 $BC,EC$ 的完备化记作 $\widehat{BC},\widehat{EC}$。$\widehat{BC},\widehat{EC}$ 是完备赋值环。
\end{definition}

语言准备完毕，进入 Gromov–Witten 势的定义。Gromov–Witten 势是同调群 $H^*(M;\C)$ 上的“函数”\footnote{之所以加引号，稍后说明。}。向量空间 $V$ 上的多项式是 $V$ 上的函数，另一方面它是 $E(V^*)$ 的元素。这里把 $V$ 的元素都看作次数 $0$。$E(V^*)$ 的完备化就是形式幂级数环，故定义 6.80 的 $\widehat{EC}$ 可看作 $C$ 的对偶上的形式幂级数。我们要定义的是 $H^*(M;\C)$ 上的“函数”，故应取 $H^*(M;\Lambda_{\mathrm{nov},\C})$ 的对偶 $C=H^*(M;\Lambda_{\mathrm{nov},\C})^*$ 上的 $\widehat{EC}$。它由 $\mathfrak{m}_k^A$ 定义。

把表示 $H_{n-*}(M;\Z)/(\operatorname{Torsion})$ 的基的几何链记作 $(P_j,\psi_j)$，$j=1,\dots,K$。（以后常省略 $\psi_j$。）Poincaré 对偶 $H_{n-*}(M;\Z)/$

% [原文 p.281]
$(\operatorname{Torsion})\simeq H^*(M;\Z)/(\operatorname{Torsion})$ 中把 $P_j$ 移过去的基，仍用同一记号表示。同调群的元素表为 $\sum z^jP_j$ 这样的和。把 $z^j\colon H^*(M;\C)\to\C$ 称为给出 $P_j$ 的系数的函数，即坐标函数。$z^j\in H^*(M;\C)^*$。

对多重下标 $J=(j_1,\dots,j_{k+1})\in\{1,\dots,b\}^{k+1}$，令 $P_{J-}=P_{j_1},\dots,P_{j_k}$、$z^J=z^{j_1}\otimes\cdots\otimes z^{j_{k+1}}\in\widehat{E}H^*(M;\C)^*$。又取 $\beta_i$，$i=0,1,2,\dots$ 使之满足条件 6.50。

\begin{definition}{6.81}
  \[
  \begin{aligned}
    \Phi_{(M,\tilde\omega_M)}
    =\sum_i\sum_{k\ge0}\frac{1}{(k+1)!}
    \sum_{J\in\{1,\dots,b\}^{k+1}}
    &\big\langle\mathfrak{m}_{k,\beta_i}^A(P_{J-}),P_{j_{k+1}}\big\rangle\\
    &\cdot T^{\beta_i\cap\tilde\omega_M}\otimes z^J .
  \end{aligned}
  \]
  这里 $\langle\cdot,\cdot\rangle$ 是相交形式 $H_k(M;\Q)\otimes H_{2n-k}(M;\Q)\to\Q$，关于 $i$ 的求和取在满足 $\deg\mathfrak{m}_{k,\beta_i}(P_{J-})+\deg P_{j_{k+1}}=2n$ 的范围内。
\end{definition}

由定义与引理 6.58，$\Phi_{(M,\tilde\omega_M)}$ 是 $\widehat{E}H^*(M;\Lambda)^*$ 的元素。$\Phi_{(M,\tilde\omega_M)}$ 的定义中含 $T$，但如前所述这是技术性理由，$\Phi_{(M,\tilde\omega_M)}$ 是 $z\in H^*(M;\C)$ 的函数（形式幂级数）。下面看 $\Phi_{(M,\tilde\omega_M)}$ 作为同调群上的函数时与 $\tilde\omega_M$ 无关。用收敛幂级数考虑时结论更易看清。因此以后用收敛幂级数考虑。明确一下定义域。以下一段时间内，对 $\Phi_{(M,\tilde\omega_M)}$ 只看 $z^J$ 中 $z^j$ 的次数为偶数的部分。记作 $\Phi^{\mathrm{even}}_{(M,\tilde\omega_M)}$。于是 $z^j$ 之间可交换，故 $z^J$ 是 $H^{\mathrm{even}}(M;\C)$ 上的多项式。

设 $\omega_0$ 是 $M$ 上的辛结构，把 $M$ 上辛结构全体之集合\footnote{指辛形式之集合，不是按辛同胚取等价类所得的集合。}中含 $\omega_0$ 的连通分支记作 $\operatorname{Sym}(M)$。以下固定连通分支 $\operatorname{Sym}(M)$，让辛结构在其中变动。把辛形式对应到它的同调类，得映射 $\mathrm{homo}\colon\operatorname{Sym}(M)\to H^2(M;\R)$。其像在乘以正常数下不变。给 $H^2(M;\R)$ 取一个范数，令
\[
  S(M)=\{[\omega]\in\mathrm{homo}(\operatorname{Sym}(M))\subset H^2(M;\R)\mid\|[\omega]\|=1\},
\]
则 $\mathrm{homo}$ 的像是 $S(M)$ 上的锥。

% [原文 p.282]
\begin{definition}{6.82}
  当把 $\Phi^{\mathrm{even}}_{(M,\tilde\omega_M)}$ 各项中的 $T$ 代入 $T=e^{-1}$ 后，所得级数在 $z\in H^{\mathrm{even}}(M;\C)$ 上绝对收敛时，记作 $\Phi^{\mathrm{even}}_{(M,\tilde\omega_M)}(z)$。
\end{definition}

设 $K$ 为 $S(M)$ 的紧子集，$K_{2i}$，$i=0,2,3,\dots$ 为 $H^{2i}(M;\C)$ 的紧子集。对正数 $S$，令 $C_{\ge S}K=\{s[\omega]\mid s\ge S,\ [\omega]\in K\}$。

\begin{conjecture}{6.83}
  对任意 $K$，存在 $S,K_i$，$i=0,2,3,\dots$，使下述成立。设 $z\in H^{\mathrm{even}}(M;\C)$ 且
  \[
    z\in C_{\ge S}K\times H^2(M;\sqrt{-1}\R)\times K_0\times K_4\times\cdots\times K_{2n},
  \]
  则 $\Phi^{\mathrm{even}}_{(M,\tilde\omega_M)}([\tilde\omega_M]+z)$ 绝对收敛。
\end{conjecture}

以下假定猜想 6.83。后面出现的命题都可以改写为关于形式幂级数的命题，这样便不必假定猜想 6.83 也能陈述结论。改写只是麻烦，本质困难没有，留给读者。以下代入都在猜想 6.83 所主张的收敛域内进行，不再一一声明。

\begin{proposition}{6.84}
  若 $\omega_1$ 与 $\omega_2$ 属于 $M$ 的辛结构集合的同一连通分支，则
  \[
    \Phi^{\mathrm{even}}_{(M,\tilde\omega_1)}(\tilde\omega_1+z)
    =\Phi^{\mathrm{even}}_{(M,\tilde\omega_2)}(\tilde\omega_2+z)
  \]
  成立。
\end{proposition}

证明中，注 6.48 与下面要说的引理 6.85 很重要。设 $(P,\psi)$ 是余维数 $2$ 的几何链，$P_1,\dots,P_k$ 任意。

\begin{lemma*}{6.85}
  \[
    \mathfrak{m}_{k+1,\beta}^A(P,P_1,\dots,P_k)
    =(\beta\cap[P])\,\mathfrak{m}_{k,\beta}^A(P_1,\dots,P_k).
  \]
\end{lemma*}

\begin{proof}
  考虑丢掉最前面那个带名点的遗忘映射
  \begin{equation}
    \mathrm{forget}_h\colon CM_{0,k+1}(M,J_M;\beta)\to CM_{0,k}(M,J_M;\beta) \tag{6.20}
  \end{equation}
  （在定理 6.59 的证明中定义过。$h(i)=i+1$。）取 $(\Sigma,\vec z,\varphi)\in CM_{0,k+1}(M,J_M;\beta)$，则 $\varphi(\Sigma)$ 的同调类是 $\beta$，故若横截正则性成立，$\varphi(\Sigma)$ 与 $\psi(P)$ 在 $\beta\cap[P]$ 个点处相交。（准确地说，连同符号在内的交点重数是 $\beta\cap[P]$。）也就是说，(6.20) 所导

  % [原文 p.283]
  出的映射 $P_\psi\times_{\mathrm{ev}_1}CM_{0,k+1}(M,J_M;\beta)\to CM_{0,k+1}(M,J_M;\beta)$ 的映射度（degree）是 $\beta\cap[P]$。与 $P_1\times\cdots\times P_k$ 取纤维积，即得引理 6.85。
\end{proof}

与引理 6.85 同样地，当 $\dim CM_{0,1}(M,J_M;\beta)=2$ 时可证\begin{equation}
  \big\langle\mathfrak{m}_{0,\beta}^A(1),P_{j_0}\big\rangle
  =(\beta\cap P_{j_0})\times\#CM_{0,0}(M,J_M;\beta) \tag{6.21}
\end{equation}
（$\#$ 是连符号一起计数的个数）。下面开始命题 6.84 的证明。

% [Tutor Note: 引理 6.85 与 (6.21) 的技巧相同：把一个“带名标记点”的模空间经遗忘映射落到少一个标记点的模空间，再与 $P$ 的 Poincaré 对偶作纤维积；纤维数正好是 $\beta\cap[P]$，于是 $\mathfrak{m}$ 上多一个 $P$ 只是乘一个 $\beta\cap[P]$。这是把上同调类的杯积翻译成计数的那一步。]

在 $H^{\mathrm{even}}(M;\C)$ 中，把 $H^2(M;\C)$ 成分特别对待，记作 $w$。即 $z'\in\bigoplus_{n\ne1}H^{2n}(M;\C)$、$z=w+z'\in H^{\mathrm{even}}(M;\C)$。设 $Q_1,\dots,Q_{b_2}$ 是 $H^2(M;\C)$ 的基。$b=b_2+b'$。令 $J=J_2+J'$，其中 $J_2$ 是对应 $H^2(M;\C)$ 成分的部分，$J'$ 是其余部分。$P_J=Q_{J_2}P_{J'}$。令 $J_2=j_{2,1}\cdots j_{2,|J_2|}$，则由引理 6.85, (6.21) 得
\begin{equation}
  \big\langle\mathfrak{m}_{k,\beta_i}^A(P_{J-}),P_{j_{k+1}}\big\rangle
  =\prod_{\ell=1}^{|J_2|}\beta_i(Q_{j_{2,\ell}})
  \big\langle\mathfrak{m}_{k,\beta_i}^A(P_{J'-}),P_{j_{k+1}}\big\rangle . \tag{6.22}
\end{equation}
这里 $J'=\emptyset$ 时，约定
\[
  \big\langle\mathfrak{m}_{k,\beta_i}^A(P_{J'-}),P_{j_{k+1}}\big\rangle
  =\#CM_{0,0}(M,J_M;\beta).
\]
把 (6.22) 代入定义 6.81，得
\[
\begin{aligned}
  \Phi^{\mathrm{even}}_{(M,\tilde\omega)}(\tilde\omega+(z'+w))
  ={}&\sum_i\sum_{\ell=0}^\infty\sum_{k=-1}^\infty\frac{1}{\ell!\,(k+1)!}
  \sum_{J'\in\{1,\dots,b'\}^{k+1}}\sum_{J_2\in\{1,\dots,b_2\}^\ell}\\
  &\big\langle\mathfrak{m}_{k+\ell-1,\beta_i}^A(Q_{J_2}P_{J'-}),P_{j_{k+1}}\big\rangle
  e^{-\beta_i\tilde\omega}(\tilde\omega+w)^{J_2}z'^{J'}\\
  ={}&\sum_i\sum_{\ell=0}^\infty\sum_{k=-1}^\infty\frac{1}{\ell!\,(k+1)!}
  \sum_{J'\in\{1,\dots,b'\}^{k+1}}\sum_{J_2\in\{1,\dots,b_2\}^\ell}\\
  &\prod_{\ell=1}^{|J_2|}\beta_i(Q_{j_{2,\ell}})
  \big\langle\mathfrak{m}_{k,\beta_i}^A(P_{J'-}),P_{j_{k+1}}\big\rangle
  e^{-\beta_i\tilde\omega}(\tilde\omega+w)^{J_2}z'^{J'}\\
  ={}&\sum_i\sum_{k=-1}^\infty\frac{e^{\beta_i\cap w}}{(k+1)!}
  \sum_{J'\in\{1,\dots,b'\}^{k+1}}
  \big\langle\mathfrak{m}_{k,\beta_i}^A(P_{J'-}),P_{j_{k+1}}\big\rangle z'^{J'} .
\end{aligned}
\]
最后一式显然与 $\tilde\omega$ 无关。\hfill$\blacksquare$

由命题 6.84，$\Phi^{\mathrm{even}}_{(M,\tilde\omega)}(\tilde\omega+z)$ 只依赖包含 $w$ 的实部的辛结构所在连通分支。下面的定义中，对 $M$ 的辛结构

% [原文 p.284]
的连通分支也一并写出。

\begin{definition}{6.86}
  设 $z=w+z'$ 是到 $H^2(M;\C)$ 成分与其余成分的分解。
  \[
  \begin{aligned}
    \Phi_M^{\mathrm{even}}(z)&=\Phi^{\mathrm{even}}_{(M,\tilde\omega)}(\tilde\omega+z)\\
    &=\sum_i\sum_{k=-1}^\infty\frac{e^{\beta_i\cap w}}{(k+1)!}
    \sum_{J'\in\{1,\dots,b'\}^{k+1}}
    \big\langle\mathfrak{m}_{k,\beta_i}^A(P_{J'-}),P_{j_{k+1}}\big\rangle z'^{J'}
  \end{aligned}
  \]
  称为 \textbf{Gromov–Witten 势}（Gromov–Witten potential）。
\end{definition}

% [Tutor Note: 定义 6.86 中 $\omega$ 的角色已并入下标：$w$ 是 $z$ 的 $H^2$ 成分，而 $\tilde\omega$ 只通过 $T=e^{-1}$ 留下 $e^{\beta_i\cap w}$ 这一项，故 $\Phi_M^{\mathrm{even}}$ 是 $H^{\mathrm{even}}(M;\C)$ 上的函数；这正是命题 6.84 的用处。]

在往下走之前，简单谈谈超流形（super manifold）。前一段只讨论了对应于偶数次同调类的 Gromov–Witten 不变量，奇数次部分当然也重要。重要的差别在于：$\deg z_{j_1},\deg z_{j_2}$ 都是奇数时，$z_{j_1}\otimes z_{j_2}=-z_{j_2}\otimes z_{j_1}$，因此由奇数次的 $z_j$ 生成的 $EH^*(M;\C)^*$ 的部分不能看作 $H^*(M;\C)$ 上的多项式或函数。用超流形这一概念，这一部分也可看作“函数”。这里不试图尽可能一般地定义超流形，只在必要时考虑\footnote{关于超流形，数学家的解说有 [67], [110], [149]。}。考虑向量丛 $V\to M$。其全空间上的函数沿纤维方向作幂级数展开，大致可看作向量丛 $E(V^*)$ 的截面全体。这里 $E(V^*)$ 是对 $V$ 的对偶 $V^*$ 逐纤维取 $E(\cdot)$ 所得之物（对称张量代数）。（但 $V^*$ 的元素都看作次数 $0$。）

现在把偶奇性的改变（parity change）沿 $E$ 的纤维方向施行，记作 $\Pi E$。即 $\Pi E$ 是其上的函数全体为 $E(\Pi V^*)$ 的截面全体的“空间”。$\Pi V^*$ 是把 $V^*$ 的元素的次数挪动 $1$ 所得。更通俗地说，$E(\Pi V^*)$ 是 $V^*$ 的外代数丛。$E(\Pi V^*)$ 的截面全体上，由外积引入积的结构，它不可交换，而是分次交换的。因此，其上函数全体为 $E(\Pi V^*)$ 截面全体的空间，作为通常的空间并不存在。其上函数全体是分次交换的空间，称为超流形。$\Pi E$ 是超流形的例

% [原文 p.285]
子\footnote{$\Pi E$ 中奇的方向、即费米子方向是 $E$ 的纤维方向，故是向量空间。（准确地说，因为把偶奇调换了，或许该叫奇向量空间。）奇的方向不是向量空间、而是弯曲的超流形起重要作用的例子，笔者不知道。}。

要用数学的方式定式化，需用层。即定义 $M$ 上的层 $\mathcal{F}$ 为 $\mathcal{F}(U)=\Gamma(U;E(\Pi V^*))$，把环化空间（ringed space）\footnote{定义见概形（scheme）的教科书，例如 [370]。} $(M,\mathcal{F})$ 看作超流形。这样看，超流形之间的映射是什么就明白了。即 $(M,\mathcal{F})$ 到 $(N,\mathcal{G})$ 的射，是 $M$ 到 $N$ 的微分同胚 $\varphi$ 与 $\varphi^*\mathcal{G}$ 到 $\mathcal{F}$ 的环层同态组成的对。

$H^*(M;\C)$ 可看作向量丛的全空间：把正交投影 $H^*(M;\C)\to H^{\mathrm{even}}(M;\C)$ 当作向量丛的投影即可。这里改变偶奇，把纤维方向、即 $H^{\mathrm{odd}}(M;\C)$ 的方向看作费米子（fermion）或奇。此时 Gromov–Witten 势就是这个超流形上的、上述意义上的函数。

下面再回到偶数部分，但若把此后所用的诸概念，例如切空间或其上的内积等，对超流形作定式化，则以下所述几乎都可原样推广到奇数次部分。不过符号要小心。

那么，用 Gromov–Witten 势来构成以 $H^{\mathrm{even}}(M;\C)$ 的开集为参数集的 Frobenius 代数族。令 $g^{jj'}=P_j\cap P_{j'}\in\Z$。次数不合时为 $0$。$g^{jj'}$ 是可逆对称 $b\times b$ 矩阵（$b$ 指偶数次 Betti 数之和）。

\begin{theorem*}{6.87}
  对 $j_1,j_2,j_3,j_4\in\{1,\dots,m\}$，下述微分方程成立。
  \begin{equation}
  \begin{aligned}
    \sum_{j,j'}g^{jj'}
    \frac{\partial^3\Phi_M^{\mathrm{even}}}{\partial z_{j_1}\partial z_{j_2}\partial z_j}
    \frac{\partial^3\Phi_M^{\mathrm{even}}}{\partial z_{j'}\partial z_{j_3}\partial z_{j_4}}
    ={}&\sum_{j,j'}g^{jj'}
    \frac{\partial^3\Phi_M^{\mathrm{even}}}{\partial z_{j_2}\partial z_{j_3}\partial z_j}
    \frac{\partial^3\Phi_M^{\mathrm{even}}}{\partial z_{j'}\partial z_{j_1}\partial z_{j_4}} .
  \end{aligned}
  \tag{6.23}
  \end{equation}
\end{theorem*}

\begin{proof}
  这是定理 6.59 的推论。即：左边在 $[\omega]+z$ 处的值中 $z^J$ 的系

  % [原文 p.286]
  数是
  \[
  \begin{aligned}
    &\sum_{J_1+J_2=J}\sum_{j,j'}g^{jj'}
    \big\langle\mathfrak{m}_{2+\#J_1}(P_{j_1},P_{j_2},P_{J_1}),P_j\big\rangle
    \big\langle\mathfrak{m}_{2+\#J_2}(P_{j'},P_{j_3},P_{J_2}),P_{j_4}\big\rangle\\
    &\quad=\sum_{J_1+J_2=J}
    \Big\langle\mathfrak{m}_{2+\#J_2}\big(\mathfrak{m}_{2+\#J_1}(P_{j_1}P_{j_2},P_{J_1}),P_{j_3},P_{J_2}\big),P_{j_4}\Big\rangle
  \end{aligned}
  \]
  同样，右边是
  \[
    \sum_{J_1+J_2=J}
    \Big\langle\mathfrak{m}_{2+\#J_2}\big(\mathfrak{m}_{2+\#J_1}(P_{j_2}P_{j_3},P_{J_1}),P_{j_1},P_{J_2}\big),P_{j_4}\Big\rangle .
  \]
  由定理 6.59，这两者相等。
\end{proof}

把 (6.23) 称为 \textbf{WDVV 方程}（Witten–Dijkgraaf–Verlinde–Verlinde equation）。

% [Tutor Note: WDVV 方程其实就是积 $\circ$ 的结合律（见下面的定义 6.88 与引理 6.89）：它把“$\mathfrak{m}_k$ 满足 $A_\infty$ 关系”这一堆无穷多条件压缩成了 $\Phi$ 的三阶导数之间的一个方程。]

取 $z_0\in H^{\mathrm{even}}(M;\C)$ 为 $\Phi_M^{\mathrm{even}}$ 收敛域上的点。把 $T_{z_0}H^{\mathrm{even}}(M;\C)$ 与 $H^{\mathrm{even}}(M;\C)$ 同一视之。$g^{jj'}$、即 Poincaré 对偶，在 $T_{z_0}H^{\mathrm{even}}(M;\C)$ 上给出复双线性非退化内积 $\langle\cdot,\cdot\rangle$。

\begin{definition}{6.88}
  用 $\langle V_1\circ V_2,V_3\rangle=V_1V_2V_3(\Phi_M^{\mathrm{even}})$ 定义 $T_{z_0}H^{\mathrm{even}}(M;\C)$ 上的积 $\circ$。右边把 $V_i$ 看作向量空间 $H^{\mathrm{even}}(M;\C)$ 上的常系数向量场。
\end{definition}

把定义 6.88 用坐标写出，即下式 (6.24)。
\begin{equation}
  \frac{\partial}{\partial z_{j_1}}\circ\frac{\partial}{\partial z_{j_2}}
  =\sum_{j,j'}g^{jj'}
  \frac{\partial^3\Phi_M^{\mathrm{even}}}{\partial z_{j_1}\partial z_{j_2}\partial z_j}
  \frac{\partial}{\partial z_{j'}} . \tag{6.24}
\end{equation}

\begin{lemma*}{6.89}
  $(T_{v_0}H^{\mathrm{even}}(M;\C),\circ,\langle\cdot,\cdot\rangle)$ 是 Frobenius 代数。
\end{lemma*}

\begin{proof}
  $\circ$ 的结合律是定理 6.87 的推论。定义 6.7 的 (7) 来自引理 6.58。其他条件显然。
\end{proof}

% [Tutor Note: 原书 p.286 引理 6.89 把切空间记作 $T_{v_0}H^{\mathrm{even}}(M;\C)$，而上文取的是基点 $z_0$，$v_0$ 应是 $z_0$ 之误，此处照抄原书。]

\begin{remark}{6.90}
  这里没有叙述，但也可以把奇数次部分包含进来。见 [222]。
\end{remark}

以上便构成了以 $H^{\mathrm{even}}(M;\C)$ 为参数的 Frobenius 代数族。这样的结构，与 Dubrovin [81] 作为 Frobenius 流形研究过的结构一致。

\begin{remark}{6.91}
  Frobenius 流形并非把 Gromov–Witten 势的性质

  % [原文 p.287]
  抽象化而发现的，它以无限可积分为起源。无限可积系与 Frobenius 流形的关系，想来在今后的 Gromov–Witten 势研究中会起中心作用，但笔者的能力不足，本书未能叙述。有兴趣的读者从 [83], [286] 附近入手，再读 [81], [82], [84], [242], [243] 等，大概就可以了。
\end{remark}

\begin{definition}{6.92}
  所谓 \textbf{Frobenius 流形}（Frobenius manifold），是指复流形 $M$，连同 $T^{1,0}M$ 上的复双线性非退化内积\footnote{$\langle\cdot,\cdot\rangle$ 不是 Hermite 度量，而是复双线性的。} $\langle\cdot,\cdot\rangle$，以及复双线性映射 $\circ\colon T_p^{1,0}M\otimes T_p^{1,0}M\to T_p^{1,0}M$，具有下述性质。
  \begin{enumerate}[label=（\arabic*）,leftmargin=2.4em]
    \item 各点的邻域内存在复坐标系 $z_i$，使
    \[
      \left\langle\frac{\partial}{\partial z_i},\frac{\partial}{\partial z_j}\right\rangle
      =\begin{cases}0,&i\ne j\ \text{时},\\ 1,&i=j\ \text{时}\end{cases}
    \]
    成立。这样的坐标 $z_i$ 称为\textbf{平坦坐标}（flat coordinate）。
    \item $\circ$ 在各切空间 $T_p^{1,0}M$ 上给出结合的 $\C$ 代数结构。
    \item 对 $T^{1,0}M$ 的截面 $V_1,V_2,V_3$，有 $\langle V_1\circ V_2,V_3\rangle=\langle V_1,V_2\circ V_3\rangle$ 成立。
    \item 对任意点，其邻域内存在函数 $\Phi$，使对平坦坐标 $z_i$ 有
    \[
      \left\langle\frac{\partial}{\partial z_i}\circ\frac{\partial}{\partial z_j},
      \frac{\partial}{\partial z_k}\right\rangle
      =\frac{\partial^3\Phi}{\partial z_i\partial z_j\partial z_k} .
    \]
    称 $\Phi$ 为\textbf{预势}（prepotential）。
  \end{enumerate}
\end{definition}

要写第 5 个条件，需要一个引理。

\begin{lemma*}{6.93}
  把 $TM^{0,1}\simeq TM$ 同一视之时，存在唯一满足下述条件的联络 $\nabla$。
  \begin{enumerate}[label=（\arabic*）,leftmargin=2.4em]
    \item $V_1\langle V_2,V_3\rangle=\langle\nabla_{V_1}V_2,V_3\rangle+\langle V_2,\nabla_{V_1}V_3\rangle$。
    \item $\nabla_{V_1}V_2-\nabla_{V_2}V_1-[V_1,V_2]=0$。
  \end{enumerate}
\end{lemma*}

% [原文 p.288]
\begin{proof}
  证明可仿照黎曼几何中 Levi-Civita 联络（Levi-Civita connection）存在唯一性的证明。进一步，由条件 (1) 知 $\nabla$ 的曲率 $[\nabla_{V_1},\nabla_{V_2}]-\nabla_{[V_1,V_2]}$ 为 $0$。
\end{proof}

（5）存在单位元 $e\in\Gamma(M;T^{1,0}M)$，使下式成立：$e\circ x=x\circ e=x,\ \nabla e=0$。

由条件 (5)，可以取 $e=\partial/\partial z_1$。此时
\[
  \frac{\partial^3\Phi}{\partial z_1\partial z_i\partial z_j}=0
\]
成立。（证明留给读者。）

\begin{theorem}{6.94}
  对辛流形 $(M,\omega_M)$，取 $\langle\cdot,\cdot\rangle$ 为 Poincaré 对偶、$\circ$ 为定义 6.88 的积、$\Phi$ 为定义 6.86 的 Gromov–Witten 势，则在 Gromov–Witten 势的收敛域上定出 Frobenius 流形的结构。
\end{theorem}

(1)～(4) 的证明由以上所述应已显然。$e$ 取 $H^0(M;\C)\simeq H_0(M;\C)$ 的生成元。(5) 由下面的引理得到。

\begin{lemma*}{6.95}
  设 $P_0\in H^0(M;\Q)$，则对任意 $P_i$（$i=1,\dots$）与 $\beta\ne0$，有 $\mathfrak{m}_{k+1,\beta}(P_0,P_1,\dots,P_k)=0$。
\end{lemma*}

\begin{proof}
  表示 $P_0$ 的几何链是 $M$ 自身（Poincaré 对偶中把同调与上同调同一视之）。因此
  \[
    \mathrm{ev}_{k+2}\colon(P_0\times\cdots\times P_k)
    \times_{\mathrm{ev}_1,\dots,\mathrm{ev}_{k+1}}\mathcal{M}_{g,k+2}(M,J_M;\beta)\to M
  \]
  经由遗忘映射 (6.25)。
  \begin{equation}
  \begin{aligned}
    (P_0\times\cdots\times P_k)&\times_{\mathrm{ev}_1,\dots,\mathrm{ev}_{k+1}}
    \mathcal{M}_{g,k+2}(M,J_M;\beta)\\
    &\longrightarrow(P_1\times\cdots\times P_k)
    \times_{\mathrm{ev}_1,\dots,\mathrm{ev}_k}\mathcal{M}_{g,k+1}(M,J_M;\beta)
  \end{aligned}
  \tag{6.25}
  \end{equation}
  (6.25) 的纤维在 $S^2$ 上具有正维数。因此，$(P_0\times\cdots\times P_k)\times_{\mathrm{ev}_1,\dots,\mathrm{ev}_{k+1}}\mathcal{M}_{g,k+2}(M,J_M;\beta)$ 在 $\mathrm{ev}_{k+2}$ 下的像的维数比虚拟维数少 $2$。由此即得引理 6.95。
\end{proof}

一般论到此为止，下面举具体例子。先算 $\C P^2$ 的情形。

% [原文 p.289]
设 $P_i\in H^{2i}(\C P^2;\Z)$，$i=0,1,2$ 为生成元。令
\begin{equation}
  N(k)=\Big\langle\mathfrak{m}_{3k-2,kP_2}\big(\underbrace{P_1,\dots,P_1}_{3k-2\ \text{个}}\big),P_1\Big\rangle . \tag{6.26}
\end{equation}
（请确认右边次数相符。）

\begin{remark}{6.96}
  $N(k)$ 是过处于一般位置的 $3k-1$ 个 $\C P^2$ 的点的 $k$ 次曲线之条数。计算这类数的问题，称为\textbf{计数几何}（enumerative geometry）的问题。
\end{remark}

由引理 6.85，
\[
  \Big\langle\mathfrak{m}_{3k-2+\ell,kP_2}
  \big(\underbrace{P_2,\dots,P_2}_{\ell\ \text{个}},
  \underbrace{P_1,\dots,P_1}_{3k-2\ \text{个}}\big),P_1\Big\rangle
  =k^\ell N(k)
\]
成立。由引理 6.95，除这些以及 $\beta=0$ 的情形外，$\mathfrak{m}_{m,\beta}$ 均为 $0$。因此，

\begin{lemma}{6.97}
  令 $v=xP_0+yP_1+zP_2$，则下式成立。
  \[
  \begin{aligned}
    \Phi^{\mathrm{even}}_{\C P^2}(v)
    &=\frac12(xy^2+x^2z)+\sum_{k,\ell}\frac{1}{(3k-1)!}\frac{1}{\ell!}
    z^{3k-1}(ky)^\ell N(k)\\
    &=\frac12(xy^2+x^2z)+\sum_k\frac{1}{(3k-1)!}z^{3k-1}e^{ky}N(k).
  \end{aligned}
  \]
\end{lemma}

令 $f(y,z)=\sum_k\frac{1}{(3k-1)!}z^{3k-1}e^{ky}N(k)$。WDVV 方程归结为下述方程。
\begin{equation}
  f_{zzz}+f_{yyy}f_{yzz}=f_{yyz}^2 . \tag{6.27}
\end{equation}
(6.27) 就是把 (6.23) 中的 $j_1,j_2,j_3,j_4$ 取为 $1,1,2,2$ 所得（注意 $g^{20}=g^{02}=g^{11}=1$，其余为 $0$）。方程 (6.27) 给出关于 $N(k)$ 的递推式。即

\begin{theorem}{6.98}
  对 $k\ge2$，
  \[
    N(k)=\sum_{a+b=k}N(a)N(b)
    \big(a^2b^2\,{}_{3k-4}C_{3a-2}-a^3b\,{}_{3k-4}C_{3a-1}\big)
  \]
  成立。
\end{theorem}

% [原文 p.290]
另一方面，$N(1)=1$。（这就是例 6.55，即量子上同调环的计算。）因此，$N(k)$ 由定理 6.98 完全决定。定理 6.98 是 Kontsevich [222] 在假定 Gromov–Witten 不变量存在的前提下给出的。$\C P^2$ 情形下 Gromov–Witten 不变量的存在性由 Ruan–Tian [319] 等人确认。

% [Tutor Note: 定理 6.98 的递推式初值 $N(1)=1$ 来自直线（$k=1$ 时 $3k-2=1$ 个点定一条直线）；$N(2)=1$（过 5 点的一条二次曲线）。这两个初值加上递推，就把全部 $N(k)$ 定死，这正是“用 WDVV 作计数”的原始范例。]

下面考虑 $M$ 是 $3$ 维且 $c_1(M)=0$ 的情形。此时 $\mathcal{M}_{0,0}(M,J_M;\beta)$ 的（虚拟）维数与 $\beta$ 无关，总为 $0$。也就是说，粗略地讲，“对每个 $\beta$，使同调类为 $\beta$ 的伪全纯映射 $\varphi\colon S^2\to M$ 只有有限个”。（当存在 $\beta=d\beta'$、$d\ge2$ 为自然数时，这一断言不成立。后文再说。）把这个数记作 $N(\beta)$。取 $H^2(M;\Z)$ 的基 $Q_1,\dots,Q_b$，取 $H^4(M;\Z)$ 的基 $P_1,\dots,P_b$ 使 $P_i\cap Q_j=\delta_{ij}$。把 $P,Q$ 分别取为 $H^0(M;\Z)$、$H^6(M;\Z)$ 的基。由次数计算与引理 6.95，一旦变量中含 $P_i$ 以外的部分，则对 $\beta\ne0$ 有 $\mathfrak{m}_{k,\beta}=0$。

进一步用引理 6.85, (6.22)，得下述引理。

\begin{lemma}{6.99}
  $\mathfrak{m}_{k,\beta}(Q_1,\dots,Q_k)=\prod_i(\beta\cap Q_i)\,N(\beta)[\beta]$。
\end{lemma}

由引理 6.99，Gromov–Witten 势便可用 $N(\beta)$ 表示。即令 $\vec z=xP+\sum z_iP_i+\sum w_iQ_i+yQ$，并把 $\beta$ 写作 $\beta=\sum\beta_iP_i$，则

\begin{lemma}{6.100}
  令 $Q_i\cap Q_j=\sum c^{ijk}P_k$，则
  \[
  \begin{aligned}
    \Phi_M^{\mathrm{even}}(\vec z)
    ={}&\Big(\frac{x^2y}{2}+\sum xv_iw_i+\frac16\sum c^{ijk}w_iw_jw_k\Big)\\
    &+\sum_\beta N(\beta)\exp\Big(\sum\beta_iw_i\Big).
  \end{aligned}
  \]
\end{lemma}

$\Phi_M^{\mathrm{even}}(\vec z)$ 的第 2 项对 $w_i$ 求导时，每一项都乘上 $\beta_i$。利用这一点可以验证：对 $\Phi_M^{\mathrm{even}}(\vec z)$，WDVV 方程自动成立。也就是说，WDVV 方程对 $3$ 维 Calabi–Yau 流形的 Gromov–Witten 不变量不带来任何新信息。

% [Tutor Note: 此处“自动成立”的机制是：$\Phi_M^{\mathrm{even}}$ 的第 2 项是 $\sum_\beta N(\beta)e^{\sum\beta_iw_i}$，它对每个 $w_i$ 的导数只是乘 $\beta_i$，于是 (6.23) 两端化为同一个关于 $\beta$ 的恒等式。因此在 3 维 Calabi–Yau 情形，WDVV 不含信息，真正的信息在 $N(\beta)$ 本身。]

这里就 $N(\beta)$ 作一点说明。利用 Calabi–Yau 流形是半正的，可以证明下面的事。（引理 6.101 在 [64] pp.208-209 中

% [原文 p.291]
是作为猜想叙述的，但本质上可用 [253] 的论证证明。）

\begin{lemma}{6.101}
  把 $J_M$ 稍作变动，则下述成立。
  \[
    \mathcal{M}_{0,0}(M,J_M;\beta)
    =\mathcal{M}_{0,0}(M,J_M;\beta)_{\mathrm{simple}}
    \cup\mathcal{M}_{0,0}(M,J_M;\beta)_{\mathrm{multi}}
  \]
  且 $\mathcal{M}_{0,0}(M,J_M;\beta)_{\mathrm{simple}}$ 是带定向的 $0$ 维流形，$\mathcal{M}_{0,0}(M,J_M;\beta)_{\mathrm{multi}}$ 的元素 $\varphi$ 是映射度 $d\ge2$ 的映射 $\pi\colon S^2\to S^2$ 与 $\mathcal{M}_{0,0}(M,J_M;\beta')_{\mathrm{simple}}$ 的元素 $\bar\varphi\colon S^2\to M$ 的合成。这里 $\beta=d\beta'$。
\end{lemma}

把 $\mathcal{M}_{0,0}(M,J_M;\beta)_{\mathrm{simple}}$ 的元素的（连符号一起数的）个数记作 $N_0(\beta)$。$\mathcal{M}_{0,0}(M,J_M;\beta)_{\mathrm{multi}}$ 部分的计入方式很微妙。映射度 $d\ge2$ 的映射 $S^2\to S^2$ 的全体具有维数，故 $\mathcal{M}_{0,0}(M,J_M;\beta)_{\mathrm{multi}}$ 的维数不为 $0$。因此 $\mathcal{M}_{0,0}(M,J_M;\beta)_{\mathrm{multi}}$ 的维数与虚拟维数不一致。这部分算作多少，答案如下。

\begin{theorem}{6.102}
  设引理 6.101 的结论对可积的 $J_M$ 成立，并且对任意满足 $d\beta'=\beta$ 的 $\beta'$，$\mathcal{M}_{0,0}(M,J_M;\beta')_{\mathrm{simple}}$ 的元素都是浸入的有理曲线\footnote{这第 2 个假定的必要性是 Tian 告诉笔者的。笔者还从 Tian 处听说，在 $\mathbb{P}^4$ 的五次超曲面的情形，$\beta$ 的次数高时大概也不成立。}，则 $N(\beta)=\sum_{d\beta'=\beta}d^{-3}N_0(\beta')$。
\end{theorem}

关于导出 $d^{-3}$ 这一系数的数学论证，见 [233], [241] 等。[64] 中也有解说（Theorem 9.2.3）。

代入引理 6.100。令 $\sum_d d^{-3}e^x=g(x)$。
\[
  \frac{d^3}{dx^3}g=\frac{e^x}{1-e^x}
\]
成立。用这个 $g$：

\begin{theorem}{6.103}
  假定与定理 6.102 相同。令 $Q_i\cap Q_j=\sum c^{ijk}P_k$，则
  \[
  \begin{aligned}
    \Phi_M^{\mathrm{even}}(\vec z)
    ={}&\Big(\frac{x^2y}{2}+\sum xv_iw_i+\frac16\sum c^{ijk}w_iw_jw_k\Big)\\
    &+\sum_\beta N_0(\beta)\,g\Big(\sum\beta_iw_i\Big).
  \end{aligned}
  \]
\end{theorem}

% [Tutor Note: 定理 6.102 的 $d^{-3}$ 来自“用 $d$ 重覆盖的参数化去数同一条有理曲线”时的重数；把它按 $g(x)=\sum_dd^{-3}e^x$ 打包，就得到定理 6.103 中分母为 $1-e^x$ 的形式，这正是 $\omega$ 退化处发散的来源。]

在本节最后，就“模空间在量子形变”稍作叙述。定义

% [原文 p.292]
6.86 的 $\Phi_M^{\mathrm{even}}(z)$ 来考虑。$z$ 的 $H^2(M;\C)$ 成分的实部曾是辛形式 $\omega$。若 $N(\beta)\ne0$，则存在实现 $\beta$ 的伪全纯曲线。因此 $\sum\beta_iv_i$ 的实部 $\beta\cap\omega$ 为正。也就是说，把定义 6.86 右边的 $\omega$ 换成 $\C\omega$ 并令 $C\to-\infty$ 时，级数越来越容易收敛。这样的极限称为\textbf{巨大体积极限}（large volume limit）。猜想 6.83 是说 Gromov–Witten 势在巨大体积极限附近收敛。离开那里、即 $\omega$ 接近 $0$ 的地方，Gromov–Witten 势被认为不收敛，而在某处发散。这被认为给出“模空间的量子形变”。也就是说，Gromov–Witten 势的定义域不是上同调群的复化、即 $\C^b$ 的开集；从巨大体积极限向内作解析延拓，它被认为会延拓为某个弯曲空间上的函数。用物理的语言说，巨大体积极限是量子效应、即定理 6.103 右边第 2 行很小的区域，是用微扰论能理解事情的区域。离开这里、进入微扰论无法理解的区域（强耦合区域）后，例如模空间甚至会因量子效应而改变拓扑。

刚才讲的是 $\omega$ 趋向 $0$ 的情形，但此外还有把 $\omega$ 动来动去、它在某处退化、即不再是辛结构的情形。考虑在 $H^{1,1}(M;\C)\subset H^2(M;\C)$ 的范围内变动的情形。能用 凯勒形式实现的 $H^{1,1}(M;\C)$ 的元素全体，称为\textbf{凯勒锥}（Kähler cone）。

先叙述 凯勒锥内点处 $\omega$ 变为奇异的情形。例如，已知 K3 曲面的 凯勒锥内点中，有对应于带奇点的 K3 曲面的点（Kobayashi–Todorov [203]）。Nakajima [277]、Bando–Kasue–Nakajima [29] 对（复 $2$ 维）Kähler–Einstein 流形的情形证明：直径有界范围内的极限是轨形。即局部可写作 $\C^2/\Gamma$，$\Gamma\subset SU(2)$。此时 $\omega$ 退化的位置，是 $M$ 中 $\C P^1$ 沿 $A,D,E$ 型 Dynkin 图配置而成的图形。换句话说，退化沿有限条有理曲线 $S^2\to M$ 发生。复 $3$ 维以上会出现未必是轨形的退化，但

% [原文 p.293]
在 $\omega$ 退化的位置，看来仍然存在有理曲线\footnote{笔者对 $3$ 维代数簇论是外行，这有多大可信度并不清楚。以下到本节末尾的叙述也一样。}。

其次考虑 凯勒锥边界附近 Gromov–Witten 势的行为。据锥定理（cone theorem，见 [263] 第 3 章），在许多情形下仍存在 $\beta\in H_2(M;\Z)$ 能由有理曲线、即全纯映射 $S^2\to M$ 实现，且 $\beta\cap\omega=0$。

如上所述，在 $\omega$ 退化处存在有理曲线时，定理 6.103 右边的第 2 行发散（分母为 $0$）。也就是说，Gromov–Witten 势发散。

这件事在物理上的解释如下。Gromov–Witten 势是把质量大的状态的效应积分掉所得的低能有效场论的关联函数；而在对应于退化 $\omega$ 的点处，质量为 $0$ 的状态会突然出现。若把这个新出现的质量 $0$ 状态的效应当作质量大的状态去积分，就会产生发散。

由 Calabi–Yau 流形作出的 Frobenius 流形，被认为在对应于退化 $\omega$ 的点处，与由另一个 Calabi–Yau 流形作出的 Frobenius 流形（或以别的方式作出的 Frobenius 流形）相连\footnote{见 [396] 等。}。把所有 Calabi–Yau 流形的模空间连起来是否会连通，这一 Reid 的空想（Reid fantasy），被认为与这样的图像有关。

\begin{remark}{6.104}
  本书对 Gromov–Witten 不变量，除亏格 $0$ 以外只触及了定义。亏格 $1$ 以上 Gromov–Witten 不变量的研究，有 Calabi–Yau 流形情形下 Bershadsky–Cecotti–Ooguri–Vafa [43] 的量子小平–Spencer 理论（quantum Kodaira–Spencer theory），以及 Fano 流形（Fano manifold，意为 $c_1(M)$ 为正）情形下 Eguchi–Hori–Xiong [87]、Eguchi–Xiong [88] 所发现的 Virasoro 代数（Virasoro algebra）的作用等。（两者都是物理学家的研究。）两者都极其有趣，但笔者能力不足，无法叙述。前者在数学立场上的研究，
\end{remark}

% 注：以下两句在原书 p.294 开头，是本条“注 6.104”的收尾；按 BRIEF §4（句子起始页在本块）写出，下一块请勿重复。
还不太多。后者在数学上的研究正兴盛起来。见 [84], [139], [147], [237] 等。

% CHUNK ch06b DONE
